\documentclass[../../main.tex]{subfiles}
\begin{document}
\begin{oframed}
  \subfile{problem}
\end{oframed}

{\bf [解答]} $k,m\ge1$と仮定すると，与式に$x=-1$を代入して$0-1=0$となり不適．故に
\begin{align}
  km=0\label{eq:1}
\end{align}
である．さらに$x=1$を代入して
\begin{align}
  2^k-1=2^m\label{eq:2}
\end{align}
である．$k=0$とすると$2^m=0$となり不適だから，\cref{eq:1}より$m=0$が従う．この時\cref{eq:2}から$k=1$となる．

以上を与式に代入して
\begin{align}
  &\dfrac{x+1}{x^l}-1=\dfrac{1}{x^n} \nonumber \\
  &(x+1)x^n-x^{(n+l)}=x^l　\nonumber \\
  &x^{(n+l)}-x^{(n+1)}-x^n+x^l=0 \label{eq:3}
\end{align}
である．($\because x\not=0$)これが$\forall x_{\not=0}$で成立するので，\cref{eq:3}は$x$の恒等式である．故に係数比較する.

まず左辺には$x^{(n+1)}$と$x^n$という次数の違う$2$つの項が存在するので，残りの$2$つの項の次数がどちらかに一致することが必要である．$n+l\ge l$だから，$n+l=n+1$，$n=l$が従い，これをといて$(n,l)=(1,1)$である．逆にこの時\cref{eq:3}は成立し，十分．

以上より，求めるのは
\begin{align}
  (k,l,m,n)=(1,1,0,1)
\end{align}
である．$\cdots$(答)

\newpage
\end{document}
